\Problem{Consulta de Percentils}

\Statement

Disposem de les dades d'un rànking de jugadors, cadascun amb el seu nom (un \verb|string|) i els seus punts (un \verb|int|). La puntuació de tots els jugadors és única (no hi ha empats). Si els ordenem per puntuació de major a menor, el \textbf{percentil} d'un jugador és el tant per cent de jugadors que té per sota (incloent-se a si mateix) en el rànking.

\medskip

Per exemple, en el següent rànking es mostren el nom, els punts i el percentil d'un grup de 5 jugadors, ordenats per punts descendentment:

\medskip 

\begin{tabular}{lcc}
    \verb|ramon| & 15 & 100\% \\
    \verb|amy| & 13 & 80\% \\
    \verb|john| & 10 & 60\% \\
    \verb|pete| & 8 & 40\% \\
    \verb|zach| & 4 & 20\% \\
\end{tabular}

\medskip

Feu un programa que llegeixi les puntuacions d'un rànking de jugadors i després contesti una sèrie de consultes, en què es demana un nom i cal contestar amb el percentil del jugador demanat. El programa s'ha de dissenyar tenint en compte que el nombre de consultes és, típicament, \emph{dos ordres de magnitud major} que el nombre de jugadors del rànking.

\Observation

En aquest problema el centre d'interès és l'eficiència. Es tracta d'aplicar els millors algorismes coneguts i trobar formes de fer el mínim nombre d'operacions per poder produir la sortida.

\medskip

L'ús de l'algorisme \verb|sort| de \verb|vector| està permès en aquest problema. Com a recordatori, per fer servir \verb|sort|, caldria:
\begin{itemize}
\item Afegir \verb|#include <algorithm>|.
\item Definir una funció de comparació \verb|fcomp| que rebi dos elements constants \verb|a| i \verb|b| i retorni \verb|true| si \verb|a| va abans que \verb|b| en la ordenació final.
\item Cridar \verb|sort(v.begin(), v.end(), fcomp)|, suposant que \verb|v| és el vector que es vol ordenar.
\end{itemize}

\Input

L'entrada té dues parts. Primer hi ha una seqüència de noms de jugadors, acabada en "\verb|---|", a on cada jugador va acompanyat de la seva puntuació. Les puntuacions de tots els jugadors són diferents, és a dir, no hi ha jugadors empatats a punts. L'ordre dels jugadors és arbitrari. Després segueix una seqüència sense sentinella de noms de jugadors dels quals es vol consultar el percentil. 

\Output

La sortida és simplement el percentil de cada jugador consultat, un per línia. Els percentils cal mostrar-los amb 2 decimals. Per mostrar dos decimals amb tots els nombres reals d'un programa se sol posar, al principi de la funció \verb|main|:

\begin{verbatim}
cout.precision(2);
cout.setf(ios::fixed);
\end{verbatim}

\SampleTwoCol
